% SGA 3, Expose VII_B, English reconstruction -- Loop 1.
% Scope: Remark 5.2.4.A in full on authority local page 94;
% stops immediately before Section 5.2.5.
% Combined-reader page: 607. Re-edition printed page: 596.
% Editorial note in scope: 161.
% Authority: Polo--Gille Exp7B-13oct24.pdf,
% SHA-256 D303690E0195654BBE7C3DAA5D89EA23E24CD709A3575217234AC19AE8836694.
% Jacob Reinhold's CC BY 4.0 Markdown was consulted as comparison lineage.
% No genuine diagram occurs in this component; all mathematics is native LaTeX.

\medskip
\noindent\textbf{注 5.2.4.A.}
\phantomsection
\label{sga3:VIIB:5.2.4.A}
\footnote{\label{sga3:VIIB:nde:161}编者增加了此注；\SGARefLink{sga3:VIIB:5.2.5}{5.2.5} 中要用到它。}
假设 \(B\) 的高度至多为 \(r+1\)，其中 \(r\in\mathbb N^*\)。则 \(I_{r+1}=I\)，并有同构
\[
\tag{1}\label{sga3:vii:remark:5-2-4-a:equation:1}
  B
  \simeq
  \frac{k[[S_1]]}{\bigl((x_1^p)\bigr)_{x_1\in S_1}}
  \widehat\otimes_k\cdots\widehat\otimes_k
  \frac{k[[S_r]]}{\bigl((x_r^{p^r})\bigr)_{x_r\in S_r}}
  \widehat\otimes_k
  \frac{k[[S_{r+1}]]}
       {\bigl((x_{r+1}^{p^{r+1}})\bigr)_{x_{r+1}\in S_{r+1}}},
\]
其中每个 \(S_n\) 为
\[
  I^2+I_{n-1}
\]
在 \(I^2+I_n\) 中的补空间。因此 \(\omega=I/I^2\) 与
\[
  \prod_{i=1}^{r+1}S_i,
\]
等同，且立即可见，对 \(n=1,\ldots,r+1\)，\(I_n\) 在 \(\omega\) 中的像 \(\omega_n\) 与
\[
  \prod_{i=1}^{n}S_i.
\]
等同。

由此得到如下推论。令
\[
  B_r=B/\check J_r,
  \qquad
  \check J_r
  =
  \bigl((x^{p^r})\bigr)_{x\in\mathfrak r(B)},
\]
并令
\[
  \mathfrak m=I/\check J_r
\]
为 \(B_r\) 的增广理想。由于 \(\check J_r\subset I^2\)，空间
\[
  \omega(r)=\mathfrak m/\mathfrak m^2
\]
与 \(\omega\) 等同。对 \(n=1,\ldots,r\)，用 \(\omega(r)_n\) 表示
\[
  \mathfrak m_n
  =
  \{\,y\in\mathfrak m\mid y^{p^n}=0\,\}.
\]
在 \(\omega(r)\) 中的像。它也是
\[
  \{\,x\in I\mid x^{p^n}\in\check J_r\,\};
\]
在 \(\omega\) 中的像；因此 \(\omega(r)_n\) 包含 \(\omega_n\)。

另一方面，同构 \textup{\SGARefLink{sga3:vii:remark:5-2-4-a:equation:1}{(1)}} 给出
\[
  \check J_r
  =
  \bigl((x_{r+1}^{p^r})\bigr)_{x_{r+1}\in S_{r+1}},
\]
从而
\[
\tag{2}\label{sga3:vii:remark:5-2-4-a:equation:2}
  B_r
  \simeq
  \frac{k[[S_1]]}{\bigl((x_1^p)\bigr)_{x_1\in S_1}}
  \widehat\otimes_k\cdots\widehat\otimes_k
  \frac{k[[S_r]]}{\bigl((x_r^{p^r})\bigr)_{x_r\in S_r}}
  \widehat\otimes_k
  \frac{k[[S_{r+1}]]}
       {\bigl((x_{r+1}^{p^r})\bigr)_{x_{r+1}\in S_{r+1}}}.
\]
由此可知，对 \(n=1,\ldots,r-1\)，\(\omega(r)_n\) 与 \(\prod_{i=1}^{n}S_i\) 等同。因此，包含关系
\[
  \omega_n\subset\omega(r)_n
\]
在 \(n=1,\ldots,r-1\) 时为等号。
